\section{Cơ sở lý thuyết}

\subsection{Thành phần của mô hình microgrid}

Theo phạm vi đề tài \cite{specs}, mô hình gồm nguồn PV, bộ điều khiển sạc, pin lưu trữ, BMS, inverter và tải. PV cung cấp điện một chiều; bộ điều khiển sạc điều chỉnh quá trình nạp theo bộ pin; BMS theo dõi trạng thái pin và cho phép hoặc ngắt các đường sạc/xả. Inverter chuyển điện DC thành AC cho tải phù hợp. Lớp IoT thu thập dữ liệu, hiển thị và tiếp nhận lệnh của người dùng.

Bộ điều khiển sạc và BMS có vai trò riêng: nhóm cần thiết kế hoặc chọn bộ sạc có đặc tuyến phù hợp, đồng thời dùng BMS để giám sát và bảo vệ bộ pin. Chức năng ngắt quá áp của BMS không thay thế việc điều khiển điện áp và dòng sạc.

\subsection{Đặc tính tấm pin quang điện}

Đường I--V biểu diễn quan hệ dòng điện và điện áp đầu ra; đường P--V được xác định từ:
\[
P(V)=V\,I(V).
\]

Các điểm đặc trưng gồm dòng ngắn mạch \(I_{\mathrm{SC}}\), điện áp hở mạch \(V_{\mathrm{OC}}\), điện áp và dòng tại điểm công suất cực đại \(V_{\mathrm{MPP}}, I_{\mathrm{MPP}}\). Công suất cực đại là:
\[
P_{\mathrm{MPP}}=V_{\mathrm{MPP}}\,I_{\mathrm{MPP}}.
\]

Nghiên cứu \cite{yaqoob} sử dụng điều kiện tham chiếu \(G_{\mathrm{ref}}=1000\ \mathrm{W/m^2}\) và nhiệt độ cell \(T_{\mathrm{ref}}=25^\circ\mathrm{C}\). Khi giữ nhiệt độ cố định, dòng quang sinh tăng gần tỷ lệ với bức xạ. Với các tấm silicon khảo sát trong bài, tăng nhiệt độ làm điện áp hở mạch giảm và công suất cực đại giảm. Nhiệt độ cell trong mô hình không được đồng nhất với nhiệt độ không khí xung quanh.

MPPT (Maximum Power Point Tracking) là quá trình điều chỉnh điểm làm việc để tìm công suất lớn. Khi có nhiều đỉnh trên đường P--V do che bóng, một phương pháp tìm kiếm cục bộ có thể dừng ở đỉnh chưa phải lớn nhất; vấn đề này liên quan đến việc tìm hiểu cấu hình dàn pin \cite{thanh}.

\subsection{Mô hình một diode}

Theo \cite{yaqoob} và \cite{pvpmc}, mô hình gồm nguồn dòng quang sinh, diode, điện trở song song và điện trở nối tiếp. Với các tham số quy đổi ở cấp module, phương trình dòng đầu ra là:
\[
I=I_{\mathrm{ph}}
-I_0\left[\exp\!\left(\frac{V+IR_s}{nN_sV_T}\right)-1\right]
-\frac{V+IR_s}{R_{\mathrm{sh}}},
\qquad
V_T=\frac{kT_K}{q}.
\]

Trong đó \(I_{\mathrm{ph}}\) là dòng quang sinh; \(I_0\) là dòng bão hòa diode; \(R_s\) và \(R_{\mathrm{sh}}\) là điện trở nối tiếp và song song; \(n\) là hệ số lý tưởng diode; \(N_s\) là số cell nối tiếp; \(k\) là hằng số Boltzmann; \(q\) là điện tích nguyên tố; \(T_K\) là nhiệt độ tuyệt đối.

Nhóm sẽ kiểm tra cách quy đổi tham số trước khi dựng mô hình. Nếu đã dùng \(R_s,R_{\mathrm{sh}}\) ở cấp module thì không nhân thêm số cell vào hai điện trở này. Tương tự, khi gộp \(nN_sV_T\) thành một hệ số, nhóm phải tránh đưa \(N_s\) vào hai lần \cite{pvpmc}.

Trong \cite{yaqoob}, tác giả biểu diễn các thành phần dòng bằng khối toán học và nguồn phụ thuộc trong Proteus. Nhóm dự kiến sử dụng bộ tham số công bố của SM55 để kiểm tra mô hình ban đầu, sau đó cập nhật theo tấm pin được lựa chọn. Các hệ số phụ thuộc nhiệt độ sẽ được đối chiếu với công thức, bảng tham số và chiều biến thiên trong tài liệu trước khi chạy.

\subsection{Che bóng một phần và tổn hao do không tương hợp}

Các module nối tiếp có cùng dòng nhánh, còn các module song song có cùng điện áp đầu cực. Khi bức xạ không đồng đều, đặc tuyến các module khác nhau nên điểm làm việc chung có thể khiến dàn pin mất công suất so với trường hợp từng module làm việc tại điểm tối ưu \cite{thanh}.

Cell bị che có thể chịu phân cực ngược và tiêu tán công suất, gây điểm nóng. Diode bypass tạo đường dòng đi vòng qua phần được bảo vệ, giúp hạn chế tác động này nhưng không khôi phục năng lượng mặt trời đã bị che. Cấu hình lại dàn pin hướng đến giảm phần tổn hao do không tương hợp thông qua thay đổi kết nối \cite{thanh}.

\subsection{Pin lưu trữ và trạng thái sạc}
\label{sec:soc}

Dung lượng điện của pin thường biểu diễn bằng Ah. Với điện áp danh định \(V_{\mathrm{nom}}\) và dung lượng \(Q_{\mathrm{Ah}}\), năng lượng danh định được ước lượng:
\[
E_{\mathrm{nom}}\approx V_{\mathrm{nom}}\,Q_{\mathrm{Ah}}\quad [\mathrm{Wh}].
\]

Đây là giá trị danh định; năng lượng cấp được cho tải còn phụ thuộc chế độ xả và tổn hao của hệ thống. C-rate biểu diễn dòng sạc/xả theo dung lượng, chẳng hạn \(1C\) tương ứng dòng có giá trị bằng dung lượng Ah chia cho một giờ \cite{kulkarni}.

SoC (State of Charge) biểu diễn tỷ lệ dung lượng còn lại. Tra điện áp hở mạch cần bảng đặc tính phù hợp và điều kiện nghỉ; coulomb counting tích phân dòng vào/ra, cần biết trạng thái ban đầu và dung lượng tham chiếu \cite{tigauge}. Nhóm đề xuất dùng coulomb counting cho mô hình chức năng, với quy ước dòng pack dương khi xả:
\[
z_{k+1}=\operatorname{clip}_{[0,1]}
\!\left(z_k-\frac{I_{\mathrm{bat},k}\,\Delta t}{3600\,Q_{\mathrm{Ah}}}\right),
\qquad
\mathrm{SoC}_k=100\,z_k.
\]

Ở công thức này, \(\Delta t\) tính bằng giây; mô hình ban đầu giả thiết hiệu suất coulomb bằng 1. Khi dòng đổi sang chiều sạc, \(I_{\mathrm{bat}}<0\) và SoC tăng. Do sai số dòng và giá trị khởi tạo gây sai số tích lũy, nhóm bổ sung hiệu chỉnh bằng OCV khi có bảng OCV--SoC và điều kiện nghỉ phù hợp \cite{tigauge}. \pending{bảng OCV--SoC và độ chính xác trên pin thật sẽ xác định sau khi chọn cell}

\subsection{Cân bằng cell, nhiệt độ và bảo vệ}

Các cell nối tiếp có thể lệch trạng thái sạc. Cell chạm giới hạn trước sẽ hạn chế khả năng sử dụng của bộ pin. Cân bằng thụ động tạo nhánh xả qua điện trở ở cell được chọn; năng lượng dư chuyển thành nhiệt. Khi áp dụng, cần xét điều kiện cho phép cân bằng, công suất tản và ảnh hưởng đến phép đo \cite{tibalance}.

Nhóm dự kiến giám sát nhiệt độ bằng LM35 hoặc mô hình tín hiệu tương đương. Theo datasheet \cite{lm35}, hệ số chuyển đổi là \(10\ \mathrm{mV}/^\circ\mathrm{C}\):
\[
T[^\circ\mathrm{C}]=\frac{V_{\mathrm{LM35}}[\mathrm{V}]}{0.01}.
\]

BMS cần phản ứng khi điện áp, dòng hoặc nhiệt độ vượt phạm vi cho phép. Các mức giới hạn phụ thuộc cell, bộ sạc và phần tử công suất được chọn. Nhóm không áp dụng trực tiếp ngưỡng của bộ pin trong \cite{kulkarni} cho một loại pin khác.

\subsection{Giám sát và điều khiển qua IoT}

Trong \cite{kulkarni}, bộ điều khiển thu thập dữ liệu cảm biến, hiển thị tại chỗ và chuyển dữ liệu qua ESP32 đến ThingSpeak. Nhóm sử dụng kiến trúc này làm cơ sở cho luồng giám sát. Các chức năng điều khiển tải, xác nhận lệnh và lưu trạng thái sự cố là phần nhóm đề xuất thêm cho hệ thống của mình.

Dữ liệu dự kiến gồm điện áp cell và pack, dòng pin, điện áp/dòng PV, công suất, nhiệt độ, SoC, trạng thái tải, nguồn đang sử dụng và mã lỗi. Thuật toán bảo vệ sẽ chạy tại bộ điều khiển cục bộ để việc ngắt mạch không phụ thuộc vào kết nối mạng.
